\documentclass{article}
\usepackage[utf8]{inputenc}
\usepackage[francais]{babel}
\usepackage{graphicx}
\usepackage[a4paper, total={6in, 8in}]{geometry}

\begin{document}

\begin{titlepage}
\newcommand{\HRule}{\rule{\linewidth}{0.5mm}} 
\center

\includegraphics[width=6cm]{assets/logo_unamur}
\\[2cm]

\textsc{\LARGE Université de Namur}
\\[2cm]

\textsc{\Large IDASM103 : Visualisation de l'information}
\\[0.2cm]

\HRule
\\[0.4cm]
\textsc{\huge Pimp My Requin}
\\[0.2cm]
\HRule
\\[0.4cm]
{\large Novembre 2023}
\\[1.5cm]

\begin{figure}[h!]
	\centering
	\includegraphics[width=10cm]{assets/logo_pimpmyrequin}
	\\[1.5cm]
\end{figure}

\begin{minipage}{0.5\textwidth}
	\begin{flushleft} 
		\emph{Auteur}
		\\
		Emeric \textsc{Lambois}
		\\
		Julien \textsc{Mernier}
		\\
		Benjamin \textsc{Pans}
		\\
		Victor \textsc{Santelé}
		\\
		Jonathan \textsc{Smith}
		\\
	\end{flushleft}
\end{minipage}
~
\begin{minipage}{0.4\textwidth}
	\begin{flushright}
		\emph{Professeur}
		\\
		Antoine \textsc{Clarinval}
		\\
	\end{flushright}
\end{minipage}

\end{titlepage}

\newpage

\tableofcontents

\newpage
\section{Abstraction des données}



\newpage
\section{Description de l’utilisateur}

Mr X est un personnage un peu particulier, depuis sa tendre enfance il fasciné par les fonds marins et plus particulièrement par les requins. Prenant un peu trop à cœur les films Jurassic Park, il a décidé de poursuivre une quête complètement folle, devenir un scientifique et de recréer un mégalodon. Mr X cherche à localiser un maximum de requin proche des plages, avec des caractéristiques essentielles pour réaliser sa nouvelle création.

\section{Abstraction des tâches}

\begin{enumerate}
	\item Choisir la partie du requin que l’on souhaite récupérer.
	\item Définir les caractéristiques de la partie du requin.
	\item Comparer les espèces de requin entre elles dans le but de définir l’espèce qui permet de récupérer plusieurs parties en une fois.
	\item Établir une liste d’espèces de requin qui constituera les étapes du voyage à réaliser pour finaliser le mégalodon.
	\item Visualiser le voyage à réaliser pour récupérer les parties du/des différents requins.
	\item Comparer la durée et le coût du trajet afin d’optimiser la récupération de plusieurs parties de requin.
\end{enumerate}

\subsection{Sélection des parties de requin}

\subsection{Ajustements des caractéristiques}

\subsection{Aide au choix des espèces}

\subsection{Sélection des espèces}

\subsection{Itinéraire du voyage}

\newpage
\section{Prototype de basse fidélité}

\begin{figure}[!h]
	\centering
	\includegraphics[width=14.4cm]{assets/prototype/basse/onglet1}
	\caption{\label{assets/prototype/basse/onglet1} Onglet 1}
\end{figure}

\vspace{0.3cm}

\begin{figure}[!h]
	\centering
	\includegraphics[width=14.4cm]{assets/prototype/basse/onglet2}
	\caption{\label{assets/prototype/basse/onglet2} Onglet 2}
\end{figure}

\begin{figure}[!h]
	\centering
	\includegraphics[width=14.4cm]{assets/prototype/basse/onglet3}
	\caption{\label{assets/prototype/basse/onglet3} Onglet 3}
\end{figure}

\vspace{1cm}

\begin{figure}[!h]
	\centering
	\includegraphics[width=14.4cm]{assets/prototype/basse/onglet4}
	\caption{\label{assets/prototype/basse/onglet4} Onglet 4}
\end{figure}

\begin{figure}[!h]
	\centering
	\includegraphics[width=14.4cm]{assets/prototype/basse/onglet5}
	\caption{\label{assets/prototype/basse/onglet5} Onglet 5}
\end{figure}

\clearpage
\section{Mapping entre les taches et le prototype}

\end{document}